\chapter{Phân Tích Và Thiết Kế Hệ Thống}

\section{Yêu cầu chức năng}

\begin{longtable}{p{1.2cm}p{5.0cm}p{8.0cm}}
\caption{Yêu cầu chức năng chính}\label{tab:functional}\\
\toprule
\textbf{Mã} & \textbf{Chức năng} & \textbf{Mô tả} \\
\midrule
\endfirsthead
\toprule
\textbf{Mã} & \textbf{Chức năng} & \textbf{Mô tả} \\
\midrule
\endhead
F01 & Chọn địa điểm & Chọn từ danh sách, tìm tên/địa chỉ hoặc nhấp bản đồ; điểm
được ánh xạ về nút graph gần nhất. \\
F02 & Thiết lập mục tiêu & Tối ưu tổng chiều dài hoặc tổng thời gian ước tính. \\
F03 & Thiết lập điểm cuối & Tuyến mở, điểm cuối cố định hoặc quay về điểm đầu. \\
F04 & Tối ưu và hiển thị & Hiển thị thứ tự ghé, polyline, khoảng cách, thời gian
và thời gian chạy thuật toán. \\
F05 & Chuyển chế độ & Chuyển giữa Lộ trình và Trốn tìm mà không mất trạng thái
của chế độ còn lại. \\
F06 & Cấu hình ván & Chọn sân, chiến thuật, seed, ngân sách, chu kỳ tiết lộ và
giới hạn lượt. \\
F07 & Chơi theo lượt & Chỉ cho chọn nút khả đạt; cập nhật tập niềm tin, nước máy
và trạng thái thắng/thua. \\
F08 & Phát lại và xuất log & Xem lại từng lượt và tải lịch sử ván ở dạng JSON. \\
\bottomrule
\end{longtable}

\section{Yêu cầu phi chức năng}

Hệ thống phải chạy cục bộ trên Python 3.11; graph định tuyến được tải một lần và
cache trong phiên; lỗi dữ liệu phải hiện thành thông báo có thể hiểu thay vì làm
sập giao diện. Phần trò chơi phải chạy ngay cả khi graph định tuyến lớn không có
sẵn. Với tối đa 10 điểm ghé, thời gian Held--Karp cần nhỏ so với bước tính ma
trận đường đi. Các trạng thái và cấu hình phải có thể tái lập bằng seed.

Về độ tin cậy, dữ liệu biên phải được kiểm tra: tọa độ hữu hạn, trọng số dương,
ngân sách dương, graph trò chơi từ 3 đến 600 nút, điểm xuất phát/đội/lối thoát
không trùng nhau và người trốn có thể tới ít nhất một lối thoát. Đây là các nghĩa
vụ triển khai tương ứng với giả định trong mô hình toán.

\section{Kiến trúc tổng thể}

\begin{figure}[H]
\centering
\begin{tikzpicture}[
  node distance=0.75cm and 0.85cm,
  box/.style={draw,rounded corners,align=center,minimum width=3.2cm,minimum height=1.0cm,fill=blue!4},
  data/.style={draw,shape=cylinder,shape border rotate=90,aspect=0.25,align=center,minimum width=2.8cm,minimum height=1.0cm,fill=teal!7},
  arr/.style={-{Latex[length=2mm]},thick}
]
\node[box] (ui) {Streamlit\\\texttt{app.py}};
\node[box,below left=of ui] (route) {Bộ điều phối tuyến\\\texttt{route\_optimizer.py}};
\node[box,below right=of ui] (chaseui) {Giao diện trò chơi\\\texttt{chase\_ui.py}};
\node[box,below=of route] (alg) {Dijkstra / Held--Karp\\NN / 2-opt};
\node[box,below=of chaseui] (game) {Luật và chính sách\\\texttt{chase.py}};
\node[data,below=of alg] (biggraph) {Graph định tuyến\\350.930 nút};
\node[data,below=of game] (smallgraph) {Sân trò chơi\\81 hoặc 358 nút};
\draw[arr] (ui) -- (route);
\draw[arr] (ui) -- (chaseui);
\draw[arr] (route) -- (alg);
\draw[arr] (chaseui) -- (game);
\draw[arr] (biggraph) -- (alg);
\draw[arr] (smallgraph) -- (game);
\end{tikzpicture}
\caption{Kiến trúc hai chế độ dùng chung giao diện nhưng tách graph theo mục đích}
\label{fig:architecture}
\end{figure}

Việc tách hai graph giải quyết hai nhu cầu khác nhau. Định tuyến cần mạng lớn và
giữ đúng chiều đường ô tô. Trò chơi cần graph nhỏ để tiền tính khoảng cách mọi
cặp và phản hồi nhanh mỗi lượt. Không dùng graph trò chơi để thay thế graph định
tuyến khi tệp lớn thiếu, vì việc đó sẽ làm thay đổi ngầm ý nghĩa ứng dụng.

\section{Luồng tối ưu lộ trình}

\begin{figure}[H]
\centering
\begin{tikzpicture}[
  node distance=0.55cm,
  flowbox/.style={draw,rounded corners,align=center,text width=10.5cm,minimum height=0.8cm,fill=blue!3},
  arr/.style={-{Latex[length=2mm]},thick}
]
\node[flowbox] (s1) {1. Chuẩn hóa điểm người dùng và tìm nút graph gần nhất};
\node[flowbox,below=of s1] (s2) {2. Chạy Dijkstra từ từng điểm quan tâm tới các điểm còn lại};
\node[flowbox,below=of s2] (s3) {3. Lập ma trận chi phí có hướng và lưu đường đi để khôi phục};
\node[flowbox,below=of s3] (s4) {4. Chạy Held--Karp theo chế độ điểm cuối; chạy NN và 2-opt để so sánh};
\node[flowbox,below=of s4] (s5) {5. Ghép các đoạn Dijkstra, tính chỉ số và vẽ tuyến trên Folium};
\draw[arr] (s1)--(s2); \draw[arr] (s2)--(s3); \draw[arr] (s3)--(s4); \draw[arr] (s4)--(s5);
\end{tikzpicture}
\caption{Pipeline tối ưu lộ trình}
\label{fig:route-pipeline}
\end{figure}

Điểm tách quan trọng là Dijkstra tối ưu từng đoạn trên graph đường, còn
Held--Karp tối ưu thứ tự trên graph đầy đủ của các điểm quan tâm. Do chi phí của
mỗi cung trong graph đầy đủ đã là đường đi ngắn nhất, ghép các cung theo thứ tự
Held--Karp tạo tuyến tối ưu đối với mô hình chi phí tĩnh và tập điểm đã ánh xạ.

\section{Thiết kế trạng thái trò chơi}

Ba kiểu dữ liệu bất biến được sử dụng:
\begin{itemize}
  \item \texttt{Rules}: ngân sách người trốn, ngân sách đội máy, chu kỳ tiết lộ,
  giới hạn lượt;
  \item \texttt{Observation}: chỉ chứa thông tin chính sách máy được phép biết;
  \item \texttt{State}: chứa vị trí thật của người trốn, quan sát công khai,
  trạng thái kết thúc và các đường đi của lượt gần nhất.
\end{itemize}

\begin{figure}[H]
\centering
\begin{tikzpicture}[
  node distance=0.7cm and 1cm,
  box/.style={draw,rounded corners,align=left,text width=5.1cm,minimum height=1.3cm},
  arr/.style={-{Latex[length=2mm]},thick}
]
\node[box,fill=red!5] (hidden) {\textbf{State riêng tư}\\vị trí thật, đường đi thật của người trốn};
\node[box,fill=green!7,right=of hidden] (obs) {\textbf{Observation công khai}\\đội máy, lối thoát, tập niềm tin, dấu vết};
\node[box,fill=blue!5,below=of obs] (policy) {\textbf{hunter\_decision}\\chỉ nhận Arena + Observation};
\node[box,fill=yellow!8,below=of hidden] (step) {\textbf{step}\\cập nhật luật, gọi policy, tạo State mới};
\draw[arr] (hidden) -- (step);
\draw[arr] (obs) -- (policy);
\draw[arr] (policy) -- (step);
\draw[arr] (step) -- (obs);
\draw[dashed,-{Latex[length=2mm]},gray] (hidden) -- node[above,sloped]{bị cấm} (policy);
\end{tikzpicture}
\caption{Ranh giới thông tin giữa trạng thái ẩn và chính sách máy}
\label{fig:information-boundary}
\end{figure}

Thiết kế này biến một yêu cầu khái niệm thành nghĩa vụ có thể kiểm thử: hai trạng
thái có cùng \texttt{Observation} nhưng vị trí người trốn khác nhau phải tạo cùng
đầu vào và cùng quyết định máy. Tuy nhiên, điều này chỉ chứng minh ranh giới tại
API chính sách hiện tại; nếu mã tương lai thêm biến toàn cục hoặc truy cập session
state bên ngoài API, bất biến phải được kiểm tra lại.

\section{Thiết kế dữ liệu}

\begin{table}[H]
\centering
\caption{Các nguồn dữ liệu của hệ thống}
\begin{tabularx}{\linewidth}{p{4.0cm}p{2.3cm}Y Y}
\toprule
\textbf{Nguồn} & \textbf{Quy mô} & \textbf{Mục đích} & \textbf{Giới hạn} \\
\midrule
\texttt{hanoi\_urban\_drive} & 350.930 nút, 864.982 cạnh & Định tuyến ô tô trên
phạm vi Hà Nội mở rộng & Tĩnh; không phản ánh ùn tắc thời gian thực \\
Lưới mô phỏng & 81 nút, 288 cạnh & Sân chuẩn, đối xứng, dễ tái lập & Không phải
đường thực tế \\
\texttt{chase\_hanoi} & 358 nút, 1.184 cạnh & Hình học phố Hà Nội cho trò chơi &
Trích từ nguồn gián đoạn, điều chỉnh hai chiều \\
\bottomrule
\end{tabularx}
\end{table}

Cạnh dài trong sân trò chơi được chia thành các đoạn không quá 90 m. Phép chia
bảo toàn tổng chiều dài và chiều cạnh; các cạnh ngược chiều dùng chung nút trung
gian khi tọa độ trùng. Mục đích là tạo đủ điểm dừng cho ngân sách 160 m/lượt,
không phải làm tăng độ chính xác địa lý.

\section{Mô hình lỗi và xử lý}

Lỗi được chia thành ba lớp. Lỗi đầu vào gồm tọa độ sai, ngân sách không dương và
điểm trùng; hệ thống từ chối sớm. Lỗi dữ liệu gồm tệp GraphML/Pickle thiếu hoặc
không hợp lệ; giao diện vẫn mở và hiển thị cảnh báo, còn chế độ trò chơi độc lập
vẫn dùng được. Lỗi kết nối chỉ ảnh hưởng geocoding, tải tile hoặc làm mới dữ liệu;
thuật toán định tuyến trên graph đóng gói không phụ thuộc tile nền.

\begin{notebox}[title={Phân biệt bản đồ nền và thuật toán}]
Tile bản đồ chỉ là ảnh nền trực quan. Một lỗi DNS hoặc chính sách tile có thể làm
nền màu xám trong khi graph, Dijkstra và tuyến kết quả vẫn hoạt động. Vì vậy kiểm
thử đúng phải kiểm tra cả giá trị thuật toán lẫn lớp hiển thị, không suy ra lỗi
thuật toán chỉ từ việc tile không tải.
\end{notebox}
